% sections/conclusion.tex -- conclusion. Every number is a \Num macro
% (generated/numbers.tex).
\section{Conclusion}\label{sec:conclusion}

We have presented a dynamic reduced basis that maintains the $\Z$ Betti numbers
of a simplicial complex exactly under insertions and deletions of simplices.
Insertion is the standard reduction step; deletion corrects every basis chain
that contains the deleted simplex, which is the rank-only form of known zigzag
and warm-start updates. Exactness holds for every order of the cells. Discrete
Morse theory enters through the order chosen at a build or rebuild: in the
collapse order every gradient pair is an apparent pair, and the static
reduction performs its column additions only in critical columns, where they are
exactly the reduction of the Morse complex. The cost of an update is
output-sensitive and cannot be bounded by the size of the update.

The experiments confirm the exactness and delimit where the method is useful.
Where the edits touch only small cycles, the dynamic reduced basis is fast and
nearly independent of the size of the complex: on the largest refinement
instance all its variants are
\NumWTwoNThirtyThreeSEightDrbSpeedupFullIdMin--\NumWTwoNThirtyThreeSEightDrbSpeedupFullIdMax{}
times faster than static recomputation, a gain that owes nothing to the Morse
order or to rebuilds. Where the cost grows with the complex, on surfaces and on
the non-locality witnesses, the average gain is at most a constant factor or
negative and the tail latency is often worse; on surfaces with known
connectivity, combinatorial counters are the better tool. A structured order
matters far more than which one: in a random order updates were up to
\NumWTwoNFifteenSEightDrbRandOverDrbIdAmort{} times slower, while the Morse order
saved nothing at build time against the insertion orders of our generators and
helped between rebuilds only on some workloads. Rebuilding paid off on cavities,
did not trigger on refinement and flag complexes, and on surfaces helped the
smaller instances but not the largest.

Open questions are maintaining the matching between rebuilds, so that the
guarantees of \cref{sec:morse} hold at all times; rebuild policies with provable
guarantees; bounds on footprints and reduction chains for structured families
of complexes and edits; other coefficients; and a comparison with
implementations of zigzag persistence.

\paragraph{Code and data.} Code, workloads, raw results and the scripts that
regenerate every number and figure are available from the author.
% TODO(author): replace with repository URL
The file \path{experiments/README.md} describes how to regenerate them, and
\path{experiments/results/WORKLOADS_MANIFEST.csv} records the generator command
and the SHA-256 of every workload instance.
